\documentclass{situationist-preprint}

\title{Linguistic Relativity, Weak and Reversible}
\msid{LNG-15-a}
\msfamily{LNG-15}
\msdiscipline{Philosophy of Language}
\mstype{Principal manuscript}
\mskeywords{linguistic relativity; Sapir--Whorf; colour categories; spatial frames of reference; number; verbal interference; conceptual schemes}

\begin{document}
\maketitle

\begin{abstract}
Does the language one speaks shape the way one thinks? Whorf held that it does, deeply: speakers of different languages dissect nature along lines laid down by their languages. The strong version of the thesis, on which language determines the concepts available to thought, is now generally rejected, but a large experimental literature reports effects of language on perception, memory and reasoning. We argue that language influences thought weakly: its effects are confined to restricted domains, chiefly categorical perception and memory for colour, spatial reasoning and exact number, and they are reversible, disappearing under verbal interference, shifting with task and context, and changing when speakers learn new categories. We define three grades of relativity, set out what each predicts, and show that the evidence from colour, space, time and number fits the weakest. This is what one should expect if language influences thought through labels and representational tools used online, not by fixing the concepts thought can employ; the view is consistent with the translatability of languages.
\end{abstract}

\section{The problem}

\citet{Sapir1929} held that human beings are at the mercy of the particular language that has become the medium of expression for their society, and \citet{Whorf1956} that we dissect nature along lines laid down by our native languages, so that speakers of languages with different grammatical categories are led to different views of the world. Whorf's examples were criticized as inferences from grammar to thought unsupported by non-linguistic evidence \citep{Pinker1994}.

The modern debate began with colour. \citet{BerlinKay1969} found that basic colour terms across languages follow constrained patterns, and \citet{Heider1972} found that speakers of Dani, a language with two basic colour terms, remember focal colours better than non-focal ones, as English speakers do.  \citet{KayKempton1984} found, however, that English speakers judged colours near the blue--green boundary as more different than Tarahumara speakers, whose language does not mark the boundary, and that the effect disappeared when the task prevented the use of the colour names. Later work, reviewed by \citet{WolffHolmes2011}, asks not whether language affects thought but how far.

We argue that the answer is: weakly. Language influences perception and memory in restricted domains, and the effects are reversible.

\section{Framework}

\begin{definition}[Grades of relativity]\label{d:grades}
Let $L$ be a language and $D$ a cognitive domain.
\begin{enumerate}
\item \emph{Strong relativity} (determinism): the concepts a speaker of $L$ can employ in $D$ are limited to those $L$ expresses.
\item \emph{Moderate relativity}: speaking $L$ produces lasting differences in the way its speakers represent $D$, which persist in non-linguistic tasks and are not removed by changes of task or by learning.
\item \emph{Weak relativity}: speaking $L$ influences perception or memory in $D$, but only in restricted domains, and the influence is reversible.
\end{enumerate}
\end{definition}

\begin{definition}[Reversibility]\label{d:rev}
An effect of language on a cognitive process is \emph{reversible} if it is eliminated or reversed (i) by interference with concurrent verbal processing, (ii) by changes in the task or context that make other cues available, or (iii) by learning or using new linguistic categories.
\end{definition}

\begin{definition}[Restricted domain]\label{d:restricted}
A domain is \emph{restricted} if it is a specific dimension of perception or representation, such as hue, spatial frame of reference or exact cardinality, and not a general capacity such as inference, conceptualization or categorization as a whole.
\end{definition}

\begin{principle}[Online mediation]\label{p:online}
If an effect of language on a task is eliminated by verbal interference, the effect is mediated by the online use of linguistic representations during the task, not by a lasting change in non-linguistic representation.
\end{principle}

Principle~\ref{p:online} is the inference experimenters themselves draw \citep{KayKempton1984,Winawer2007}: verbal interference, such as rehearsing digits, occupies the resources used to produce and maintain words, and an effect that disappears under it but survives equally demanding spatial interference depends on the words.

\section{The evidence}

\begin{lemma}[Colour]\label{l:colour}
The effects of language on colour cognition are restricted and reversible.
\end{lemma}

\begin{proof}
\citet{Winawer2007} found that Russian speakers, whose language obligatorily distinguishes lighter blues (\emph{goluboy}) from darker blues (\emph{siniy}), discriminate colours across that boundary faster than within it, while English speakers show no such advantage; the advantage disappeared under verbal interference but not under spatial interference. \citet{Gilbert2006} found that the categorical advantage for colours across the blue--green boundary is present when targets appear in the right visual field, which projects to the left, language-dominant hemisphere, and absent in the left visual field, and that verbal interference removes it. \citet{Athanasopoulos2009} found that Greek speakers, whose language distinguishes light and dark blue, show a reduced sensitivity to that distinction the longer they have lived in an English-speaking environment. The domain is a single dimension of perception, so it is restricted (Definition~\ref{d:restricted}); the effects are removed by verbal interference and shift with exposure to another language, so they are reversible under clauses (i) and (iii) of Definition~\ref{d:rev}.
\end{proof}

Some effects appear early in processing: \citet{Thierry2009} found that the brain responses of Greek speakers to deviant colours in an unattended stream distinguish light and dark blue, while those of English speakers do not. The influence of language thus reaches early perception; but the result of \citet{Athanasopoulos2009} on the same contrast shows that it is reversible there too, which suggests that the visual system adapts to the categories a speaker habitually uses rather than being fixed by them. \citet{Roberson2000} found differences between Berinmo and English speakers in categorical perception and memory, but not an inability to discriminate unnamed colours.

\begin{lemma}[Space]\label{l:space}
The effects of language on spatial reasoning are restricted and reversible.
\end{lemma}

\begin{proof}
\citet{Levinson2003} reported that speakers of languages that use absolute frames of reference (\enquote{north of}, \enquote{uphill}) rather than relative ones (\enquote{left of}) tend to reconstruct arrays in absolute terms after being rotated, while speakers of Dutch and English reconstruct them in relative terms. \citet{LiGleitman2002} found that English speakers' choice of frame shifts with the availability of environmental landmarks: when tested outdoors or with salient landmarks, they too reconstruct arrays in absolute terms. The effect concerns the choice of a frame of reference in an ambiguous task, a restricted domain, and it changes with the cues the task makes available, so it is reversible under clause (ii) of Definition~\ref{d:rev}.
\end{proof}

\begin{lemma}[Number]\label{l:number}
The effects of language on numerical cognition are restricted, and they are the effects of the presence or absence of a representational tool.
\end{lemma}

\begin{proof}
Speakers of Pirahã, a language without exact number words, perform accurately on tasks that require exact matching of sets when the sets remain visible, and inaccurately when the task requires holding an exact large quantity in memory \citep{Gordon2004,Frank2008}. The domain is exact cardinality beyond small numbers, which is restricted. \citet{Frank2008} interpret the result as showing that number words are a cognitive technology for remembering exact quantities, not a condition of the concept of exact number, since the one-to-one matching that the concept requires is available to the speakers when memory is not taxed. What the language lacks is a tool, and what a tool provides can be acquired with the tool, as children acquiring a counting system do.
\end{proof}

\section{The thesis}

\begin{theorem}\label{t:main}
Language influences thought weakly: it influences perception and memory in restricted domains, and the effects are reversible. Neither strong nor moderate relativity is supported.
\end{theorem}

\begin{proof}
By Lemmas~\ref{l:colour}--\ref{l:number}, the established effects concern colour discrimination and memory, choice of spatial frame and memory for exact quantity, and the reported effects on time, discussed below, concern temporal priming: each a restricted domain by Definition~\ref{d:restricted}. By Lemmas~\ref{l:colour} and~\ref{l:space}, the effects in the first two are reversible by verbal interference, by change of task cues or by learning. The reported effect of Mandarin vertical time metaphors on temporal priming \citep{Boroditsky2001} failed to replicate \citep{Chen2007,JanuaryKako2007}, and English speakers briefly trained on vertical terms showed it in the original study, so it is at most reversible under clause (iii). For number, Lemma~\ref{l:number} shows that the effect is the absence of a tool for memory, removable by acquiring the tool, which is reversibility under clause (iii). So the evidence supports weak relativity (Definition~\ref{d:grades}). Strong relativity predicts that speakers cannot employ concepts their language does not express; in every domain examined, speakers discriminate, match or reason with distinctions their language does not mark, so the prediction fails. Moderate relativity predicts effects that persist in non-linguistic tasks and resist change; the interference, cue and learning results show that they do not.
\end{proof}

The pattern has an explanation. By Principle~\ref{p:online}, the effects that disappear under verbal interference depend on the online use of words during the task. Words serve as labels that sharpen category boundaries, as compact codes for memory and, for number, as a system for representing quantities exactly: ways in which language is used in thinking, not ways in which it fixes what can be thought. The effects are restricted to domains in which a linguistic code helps with the task, and reversible because what is used online can be blocked, replaced by other cues, or relearned.

\section{Objections and replies}

\begin{objection}[Thinking for speaking]\label{o:habitual}
The interesting claim was never that language limits what can be thought, but that it shapes what speakers habitually attend to. Reversibility under laboratory manipulations is consistent with large habitual effects in everyday life, and so the weak thesis understates the influence of language.
\end{objection}

\begin{reply}
We accept that language shapes habitual attention in these domains; that is part of weak relativity. The thesis concerns depth: an effect that disappears when the words are not available, and changes when speakers learn other words, is an influence exercised through the words, not a reorganization of thought that would persist without them.
\end{reply}

\begin{objection}[Conceptual schemes]\label{o:schemes}
If languages influence thought at all, speakers of different languages inhabit different conceptual schemes, and translation between them is at best partial.
\end{objection}

\begin{reply}
The effects established are differences in speed, accuracy and default choices in restricted domains. They do not make any sentence of one language untranslatable into another. As LNG-06 argues, translation is constrained by more than stimulus meaning, and the constraints suffice to select translations; \citet{Davidson1974} argued that the idea of a conceptual scheme untranslatable into ours is incoherent.
\end{reply}

\begin{objection}[Late acquisition]\label{o:late}
Adults whose language lacks number words may find a counting system hard to acquire. An effect that cannot in practice be undone is not reversible.
\end{objection}

\begin{reply}
Difficulty in acquiring a tool late shows that tools are easier to learn early, not that the capacity the tool serves is absent. The concept of exact equality is available to such speakers, who match sets exactly when memory is not taxed (Lemma~\ref{l:number}); what is missing is the means of remembering exact quantities, and children who acquire that means acquire the capacity with it.
\end{reply}

\section{What the result does not establish}

The argument relies on domains that have been studied experimentally. It does not exclude effects in domains not yet studied, though it predicts that any such effect will be restricted and reversible if it is mediated by the online use of words. It does not concern the role of language in the development of thought in children, which may be greater. Nor does it bear on the reality of the properties that languages categorize differently: on the view of MET-17, colours are real relational properties, which languages divide differently.

\section{Conclusion}

Language influences thought weakly. The effects established in colour, space, time and number are confined to specific dimensions of perception and memory, and they are reversible: they disappear under verbal interference, shift with the cues a task provides, and change when speakers learn new categories. They are what one would expect if words are used as tools in thinking, and they give no support to the view that language fixes the concepts available to thought.

\begin{thebibliography}{99}\small
\bibitem[Athanasopoulos(2009)]{Athanasopoulos2009} Athanasopoulos, P. (2009). Cognitive representation of colour in bilinguals: The case of Greek blues. \emph{Bilingualism: Language and Cognition}, 12(1), 83--95.
\bibitem[Berlin and Kay(1969)]{BerlinKay1969} Berlin, B. and Kay, P. (1969). \emph{Basic Color Terms: Their Universality and Evolution}. University of California Press.
\bibitem[Boroditsky(2001)]{Boroditsky2001} Boroditsky, L. (2001). Does language shape thought? Mandarin and English speakers' conceptions of time. \emph{Cognitive Psychology}, 43(1), 1--22.
\bibitem[Chen(2007)]{Chen2007} Chen, J.-Y. (2007). Do Chinese and English speakers think about time differently? Failure of replicating Boroditsky (2001). \emph{Cognition}, 104(2), 427--436.
\bibitem[Davidson(1974)]{Davidson1974} Davidson, D. (1974). On the very idea of a conceptual scheme. \emph{Proceedings and Addresses of the American Philosophical Association}, 47, 5--20.
\bibitem[Frank et~al.(2008)]{Frank2008} Frank, M.~C., Everett, D.~L., Fedorenko, E. and Gibson, E. (2008). Number as a cognitive technology: Evidence from Pirah\~a language and cognition. \emph{Cognition}, 108(3), 819--824.
\bibitem[Gilbert et~al.(2006)]{Gilbert2006} Gilbert, A.~L., Regier, T., Kay, P. and Ivry, R.~B. (2006). Whorf hypothesis is supported in the right visual field but not the left. \emph{Proceedings of the National Academy of Sciences}, 103(2), 489--494.
\bibitem[Gordon(2004)]{Gordon2004} Gordon, P. (2004). Numerical cognition without words: Evidence from Amazonia. \emph{Science}, 306(5695), 496--499.
\bibitem[Heider(1972)]{Heider1972} Heider, E.~R. (1972). Universals in color naming and memory. \emph{Journal of Experimental Psychology}, 93(1), 10--20.
\bibitem[January and Kako(2007)]{JanuaryKako2007} January, D. and Kako, E. (2007). Re-evaluating evidence for linguistic relativity: Reply to Boroditsky (2001). \emph{Cognition}, 104(2), 417--426.
\bibitem[Kay and Kempton(1984)]{KayKempton1984} Kay, P. and Kempton, W. (1984). What is the Sapir--Whorf hypothesis? \emph{American Anthropologist}, 86(1), 65--79.
\bibitem[Levinson(2003)]{Levinson2003} Levinson, S.~C. (2003). \emph{Space in Language and Cognition: Explorations in Cognitive Diversity}. Cambridge University Press.
\bibitem[Li and Gleitman(2002)]{LiGleitman2002} Li, P. and Gleitman, L. (2002). Turning the tables: Language and spatial reasoning. \emph{Cognition}, 83(3), 265--294.
\bibitem[Pinker(1994)]{Pinker1994} Pinker, S. (1994). \emph{The Language Instinct}. William Morrow.
\bibitem[Roberson et~al.(2000)]{Roberson2000} Roberson, D., Davies, I. and Davidoff, J. (2000). Color categories are not universal: Replications and new evidence from a stone-age culture. \emph{Journal of Experimental Psychology: General}, 129(3), 369--398.
\bibitem[Sapir(1929)]{Sapir1929} Sapir, E. (1929). The status of linguistics as a science. \emph{Language}, 5(4), 207--214.
\bibitem[Thierry et~al.(2009)]{Thierry2009} Thierry, G., Athanasopoulos, P., Wiggett, A., Dering, B. and Kuipers, J.-R. (2009). Unconscious effects of language-specific terminology on preattentive color perception. \emph{Proceedings of the National Academy of Sciences}, 106(11), 4567--4570.
\bibitem[Whorf(1956)]{Whorf1956} Whorf, B.~L. (1956). \emph{Language, Thought, and Reality: Selected Writings of Benjamin Lee Whorf}, ed. J.~B. Carroll. MIT Press.
\bibitem[Winawer et~al.(2007)]{Winawer2007} Winawer, J., Witthoft, N., Frank, M.~C., Wu, L., Wade, A.~R. and Boroditsky, L. (2007). Russian blues reveal effects of language on color discrimination. \emph{Proceedings of the National Academy of Sciences}, 104(19), 7780--7785.
\bibitem[Wolff and Holmes(2011)]{WolffHolmes2011} Wolff, P. and Holmes, K.~J. (2011). Linguistic relativity. \emph{Wiley Interdisciplinary Reviews: Cognitive Science}, 2(3), 253--265.
\end{thebibliography}

\end{document}
